\documentclass[conference]{IEEEtran}
\usepackage{cite}
\usepackage{amsmath,amssymb,amsfonts}
\usepackage{booktabs}
\usepackage{graphicx}
\usepackage{url}
\hyphenation{op-tical net-works semi-conduc-tor}

\begin{document}

\title{Fault Aware Training For Spiking Neural Networks}

\author{
\IEEEauthorblockN{Pin-Yu Wu}
\IEEEauthorblockA{National Taiwan University\\
Taipei, Taiwan\\
b11502041@ntu.edu.tw}
\and
\IEEEauthorblockN{Yu-Shiang Wang}
\IEEEauthorblockA{National Taiwan University\\
Taipei, Taiwan\\
b11504024@ntu.edu.tw}
}

\maketitle

\begin{abstract}
Spiking Neural Networks (SNNs) are attractive for neuromorphic hardware
because they are event-driven and energy-efficient. That hardware, however,
is vulnerable to neuron and synapse faults. This paper describes a Fault
Aware Training (FAT) framework for SNNs: faults are injected in the
training forward pass and a fresh fault list is drawn every mini-batch, so
healthy neurons learn to compensate. After each optimiser step the incoming
weights of the faulty neurons are restored (weight recovery) so the network
does not overfit a particular fault pattern. On MNIST, FAT matches
fault-free accuracy of ordinary back-propagation and is more robust under
injected hardware faults. The implementation, including a normal-BP
baseline and a Table~I reproduction script, is released with this paper.
\end{abstract}

\begin{IEEEkeywords}
Spiking Neural Networks, Fault Tolerance, Neuromorphic Computation,
Robust Training, Hardware Faults.
\end{IEEEkeywords}

\section{Introduction}

Spiking Neural Networks process information with discrete spikes and are a
natural fit for neuromorphic chips. Once the network is mapped onto
silicon, manufacturing defects, ageing and environmental noise can silence
neurons, lock them on, or freeze synaptic weights. Classical hardware
remedies (redundancy, ECC) spend the very energy SNNs are meant to save.

Software-level fault tolerance trains the network so that it keeps working
when some units fail. We follow Zahid \textit{et al.}~\cite{zahid2020fat},
who trained conventional networks under simulated hardware faults, and
adapt the idea to an SNN. The goal is image recognition that cannot afford
silent failures: the network should stay accurate both when the chip is
healthy and when a random fraction of neurons is broken.

\section{Background}

Unlike a ReLU network, an SNN is not differentiable at the spike. We use
an arctan surrogate gradient~\cite{eshraghian2023snn} and clip the
gradient norm to keep training stable. Related SNN fault-tolerance work
includes neuron-level analysis~\cite{spyrou2021nft} and memory-aware
schemes such as ReSpawn~\cite{putra2021respawn}. Our contribution is a
training recipe, not a hardware encoder: three fault models, per-batch
resampling, and an optional weight-recovery step that freezes the
incoming weights of the units that were faulty in that batch.

\section{Proposed Technique}

The network is fully connected, $784 \rightarrow 256 \rightarrow 32
\rightarrow 10$, with no biases. Each linear map is followed by a leaky
integrate-and-fire (LIF) neuron (snntorch \texttt{Leaky}, $\beta=0.9$,
threshold $0.5$, subtract reset). MNIST pixels in $[0,1]$ are treated as
spike probabilities: at each of $T=25$ time steps a Bernoulli spike train
is drawn and pushed through the three stages; output spikes are summed
and used as logits. \textbf{Training, validation and test share this
forward path.}

\subsection{Fault simulation}
Every \emph{hidden} LIF neuron is declared faulty independently with
probability $p$ (default $0.25$). The class (output) layer is left
intact: at $p=0.25$ a ten-way readout would otherwise lose two or three
digits at random and the test accuracy would jump by tens of points.
Three models are supported:

\begin{itemize}
  \item \textbf{Hard-to-spike (hts)} --- the neuron never emits a spike.
  \item \textbf{Saturated (sat)} --- the neuron spikes at every time step.
  \item \textbf{Stuck-at weight (swf)} --- one random incoming synapse is
        frozen at a value drawn uniformly from $[-10,10]$.
\end{itemize}

A new list is drawn every mini-batch so the network cannot specialise to
one broken subset. The default Table~I experiment uses hts.

\subsection{Weight recovery}
After the Adam step, the incoming weight rows of the neurons that were
faulty in that batch are written back to their pre-step values. Those
units therefore do not learn from a batch in which they were broken;
healthy units absorb the update. The step can be disabled
(\texttt{-{}-no-recovery}) for ablation.

\subsection{Gradient clipping}
The gradient $\ell_2$ norm is clipped to $10$ before the optimiser step.

\subsection{Evaluation}
Two checkpoints are written each run: the best fault-free validation
accuracy and the best \emph{faulty} validation accuracy. FAT early-stops
and reports the latter; the normal-BP baseline reports the former.
Test numbers use the full $10\,000$-image MNIST set. Fault-free accuracy
is repeated over $R$ independent rate-coding draws. Faulty accuracy draws
$R$ independent fault lists; each list is applied to the \emph{whole}
test set (one draw $=$ one simulated chip). We report mean $\pm$
standard deviation over the $R$ draws, not per-batch accuracies.

\texttt{good\_inference()} is the shared forward pass with an empty
fault list; \texttt{bad\_inference()} injects a caller-supplied list.
Ordinary back-propagation (\texttt{-{}-mode normal}) skips fault
injection and weight recovery and is the Table~I baseline.

\section{Experiments}

Both methods use Adam ($3\times 10^{-4}$, weight decay $10^{-5}$), batch
size $64$, an $8{:}2$ train/validation split of MNIST, early stopping
with patience $5$, and seed-controlled data, spikes and fault draws.
Table~I is produced by \texttt{bash run\_paper\_experiments.sh}
(10 seeds $\times$ \{normal, FAT\}, 20 epochs, $R=10$ draws, hts,
output layer healthy). A two-epoch CPU sketch that still faulted the
output layer is archived in \texttt{docs/results/} and must not be mixed
with Table~I. An earlier manuscript draft that reported $\approx 82\%$
and a $2$\,pp gain used analog training inputs and per-batch test
scores; those numbers are obsolete.

\begin{table}[t]
\caption{Accuracy (\%) of normal back-propagation vs.\ FAT under
hidden-layer hard-to-spike faults, $p=0.25$. Fill from
\texttt{outputs/paper/table1/table1.md} after
\texttt{bash run\_paper\_experiments.sh}.}
\label{tab:table1}
\centering
\begin{tabular}{crrrr}
\toprule
seed & N ff & N faulty & FAT ff & FAT faulty \\
\midrule
\multicolumn{5}{l}{\textit{see outputs/paper/table1/table1.md}} \\
\bottomrule
\end{tabular}
\end{table}

Ablations (\texttt{run\_ablations.py}) train under hts / sat / swf, with
and without weight recovery and per-batch resampling, then test each
model on all three fault types. Optional extras: Fashion-MNIST
(\texttt{-{}-dataset fashion}) and a one-convolution SNN
(\texttt{-{}-arch conv}). A paired $t$-test on the per-seed differences
accompanies the mean improvement.

\section{Conclusion}

FAT injects hardware faults into the SNN forward pass, resamples them
every batch, and optionally freezes the incoming weights of the units
that failed. The same multi-step rate-coded path is used at train and
test time. The recipe is a software complement to neuromorphic fault
tolerance and is fully reproducible from the accompanying code.

\begin{thebibliography}{00}
\bibitem{zahid2020fat}
U.~Zahid, G.~Gambardella, N.~J. Fraser, M.~Blott and K.~Vissers,
``FAT: Training Neural Networks for Reliable Inference Under Hardware Faults,''
in \emph{Proc.\ IEEE ITC}, 2020, pp.~1--10.
\bibitem{eshraghian2023snn}
J.~K. Eshraghian \textit{et al.},
``Training Spiking Neural Networks Using Lessons From Deep Learning,''
\emph{Proc.\ IEEE}, vol.~111, no.~9, pp.~1479--1495, Sep.~2023.
\bibitem{spyrou2021nft}
T.~Spyrou \textit{et al.},
``Neuron Fault Tolerance in Spiking Neural Networks,''
in \emph{Proc.\ DATE}, 2021, pp.~743--748.
\bibitem{putra2021respawn}
R.~V.~W. Putra, M.~A. Hanif and M.~Shafique,
``ReSpawn: Energy-Efficient Fault-Tolerance for Spiking Neural Networks considering Unreliable Memories,''
in \emph{Proc.\ ICCAD}, 2021, pp.~1--9.
\end{thebibliography}

\end{document}
